\documentclass[11pt,a4paper]{article}

\usepackage[utf8]{inputenc}
\usepackage[T1]{fontenc}
\usepackage{amsmath,amssymb}
\usepackage{booktabs}
\usepackage{geometry}
\usepackage{algorithm}
\usepackage{algpseudocode}
\usepackage{hyperref}
\usepackage{natbib}

\geometry{margin=2.5cm}

\title{Methods: Kinetic Segregation Model of TCR--CD45 Spatial
  Exclusion on a 2D Membrane Patch}
\author{Based on Neve-Oz, Sherman \& Raveh,
  \textit{Frontiers in Immunology}, 2024}
\date{}

\begin{document}
\maketitle

%% ----------------------------------------------------------------
\section{Model overview}

The kinetic segregation (KS) model simulates the spatial segregation
of T-cell receptor (TCR) and CD45 molecules on a
$2\;\mu\text{m}\times2\;\mu\text{m}$ membrane patch.  A 2D membrane
height field $h(x,y)$ is discretised on a regular grid of spacing
$\Delta x = L / N_\text{grid}$, where $L=2000$~nm and
$N_\text{grid}$ is typically 64.  Static pMHC ligands
($N_\text{pMHC}$, default 200) are placed on the opposing
antigen-presenting cell (APC) surface and gate TCR--pMHC binding:
TCR molecules ($N_\text{TCR}$, default 50) are attracted to
tight-contact regions (low $h$) \emph{only in grid cells that
contain at least one pMHC}.  CD45 molecules ($N_\text{CD45}$,
default 100) are sterically excluded from regions where $h$ is
smaller than the CD45 ectodomain height
$h_\text{CD45}=50$~nm.

The system is evolved with a Metropolis--Hastings Monte~Carlo (MC)
algorithm at thermal energy $k_BT=1$ (all energies are expressed in
units of $k_BT$).


%% ----------------------------------------------------------------
\section{System energy}

The total energy of the system is the sum of three contributions:

\begin{equation}
  \label{eq:total}
  E_\text{total}
  \;=\;
  E_\text{bend}\!\bigl[h\bigr]
  \;+\;
  \sum_{i=1}^{N_\text{TCR}}
    \mathbb{1}_\text{pMHC}(\mathbf{r}_i)\;
    E_\text{TCR}\!\bigl(h(\mathbf{r}_i)\bigr)
  \;+\;
  \sum_{j=1}^{N_\text{CD45}} E_\text{CD45}\!\bigl(h(\mathbf{r}_j)\bigr)
\end{equation}

\noindent
where $\mathbf{r}_i$ and $\mathbf{r}_j$ are the lateral positions of
individual TCR and CD45 molecules, respectively,
$h(\mathbf{r})$ is the membrane height at that position (obtained by
nearest-grid-point interpolation), and
$\mathbb{1}_\text{pMHC}(\mathbf{r})$ is unity when the grid cell at
$\mathbf{r}$ contains at least one pMHC and zero otherwise.


%% ----------------------------------------------------------------
\section{Potential energy terms}
\label{sec:potentials}

Table~\ref{tab:potentials} summarises all energy contributions, their
functional forms, and default parameter values.

\begin{table}[ht]
  \centering
  \caption{Energy terms in the kinetic segregation model.  All
    energies are in units of $k_BT$; lengths are in nm.}
  \label{tab:potentials}
  \small
  \begin{tabular}{@{}p{2.8cm}p{5.2cm}p{3.5cm}p{3.5cm}@{}}
    \toprule
    \textbf{Term} & \textbf{Expression} & \textbf{Parameters} &
      \textbf{Description} \\
    \midrule
    Membrane bending &
      $\displaystyle E_\text{bend}
        = \frac{\kappa}{2}\sum_{i,j}
          \bigl(\nabla^2_d h_{ij}\bigr)^2\,\Delta x^2$ &
      $\kappa$: bending rigidity ($k_BT$); \newline
      $\Delta x = L/N_\text{grid}$ &
      Helfrich bending energy on a discrete grid with periodic BCs \\[6pt]
    TCR--pMHC binding &
      $\displaystyle E_\text{TCR}(h)
        = -U_\text{assoc}\,
          \exp\!\Bigl(-\frac{h^2}{2\sigma_\text{bind}^2}\Bigr)$ &
      $U_\text{assoc}=20\;k_BT$; \newline
      $\sigma_\text{bind}=3$~nm &
      Gaussian attractive well centred at tight contact ($h=0$) \\[6pt]
    CD45 steric repulsion &
      $\displaystyle E_\text{CD45}(h)
        = \begin{cases}
            \tfrac{1}{2}\,k_\text{rep}\,(h_\text{CD45}-h)^2
              & h < h_\text{CD45} \\
            0 & h \ge h_\text{CD45}
          \end{cases}$ &
      $k_\text{rep}=1.0\;k_BT/\text{nm}^2$; \newline
      $h_\text{CD45}=50$~nm &
      Harmonic repulsive wall; excludes CD45 from tight contact \\
    \bottomrule
  \end{tabular}
\end{table}


\subsection{Membrane bending energy}

The bending energy follows the Helfrich Hamiltonian discretised on a
square grid with periodic boundary conditions.  The discrete Laplacian
at grid point $(i,j)$ is computed via the standard five-point stencil:

\begin{equation}
  \label{eq:laplacian}
  \nabla^2_d\, h_{ij}
  \;=\;
  \frac{h_{i-1,j} + h_{i+1,j} + h_{i,j-1} + h_{i,j+1} - 4\,h_{ij}}
       {\Delta x^2}
\end{equation}

\noindent
with indices taken modulo $N_\text{grid}$ (periodic boundaries).  The
full bending energy is then:

\begin{equation}
  \label{eq:bending}
  E_\text{bend}
  \;=\;
  \frac{\kappa}{2}
  \sum_{i=0}^{N_\text{grid}-1}\;
  \sum_{j=0}^{N_\text{grid}-1}
    \bigl(\nabla^2_d\, h_{ij}\bigr)^2\;\Delta x^2
\end{equation}


\subsection{TCR--pMHC binding potential}

TCR molecules are attracted to regions of low membrane height
(tight contact between the T-cell and antigen-presenting cell),
modelling TCR--pMHC engagement.  The potential is a Gaussian well
centred at $h=0$:

\begin{equation}
  \label{eq:tcr}
  E_\text{TCR}(h) \;=\; -U_\text{assoc}\,
    \exp\!\left(-\frac{h^2}{2\,\sigma_\text{bind}^2}\right)
\end{equation}

\noindent
where $U_\text{assoc}=20\;k_BT$ is the association energy depth and
$\sigma_\text{bind}=3$~nm is the spatial width of the binding well.


\subsection{CD45 steric repulsion}

CD45 has a large extracellular domain ($\sim$35~nm) and is sterically
excluded from tight-contact zones.  This is modelled as a one-sided
harmonic potential:

\begin{equation}
  \label{eq:cd45}
  E_\text{CD45}(h) \;=\;
    \begin{cases}
      \displaystyle\frac{1}{2}\,k_\text{rep}\,
        \bigl(h_\text{CD45} - h\bigr)^2
        & \text{if } h < h_\text{CD45} \\[6pt]
      0 & \text{otherwise}
    \end{cases}
\end{equation}

\noindent
with $k_\text{rep} = 1.0\;k_BT/\text{nm}^2$ and
$h_\text{CD45} = 50$~nm.


%% ----------------------------------------------------------------
\section{Monte Carlo algorithm}
\label{sec:mcmc}

The simulation uses a Metropolis--Hastings MC scheme with two
alternating update phases per sweep.  Each \emph{sweep} updates every
degree of freedom once (all molecules and all grid cells), ensuring
ergodic sampling of the configuration space.

Algorithm~\ref{alg:mc} gives the complete pseudocode.

\begin{algorithm}[ht]
\caption{Metropolis Monte Carlo for kinetic segregation}
\label{alg:mc}
\begin{algorithmic}[1]
\Require $N_\text{steps}$, membrane grid $h$, positions
  $\{\mathbf{r}^\text{TCR}_i\}$, $\{\mathbf{r}^\text{CD45}_j\}$,
  parameters $\kappa, U_\text{assoc}, \sigma_\text{bind},
  h_\text{CD45}, k_\text{rep}$
\Ensure Final configuration $\{h, \mathbf{r}^\text{TCR},
  \mathbf{r}^\text{CD45}\}$
\Statex
\For{$s = 1$ \textbf{to} $N_\text{steps}$}
  \Statex \Comment{\textbf{Phase 1: Molecular moves}}
  \For{each molecule $k$ in TCR $\cup$ CD45}
    \State $\mathbf{r}_\text{old} \gets \mathbf{r}_k$
    \State $E_\text{old} \gets E_\text{mol}(\mathbf{r}_k, h)$
      \Comment{Eq.~\ref{eq:tcr} or \ref{eq:cd45}}
    \State $\mathbf{r}_k \gets \mathbf{r}_k
      + \mathcal{N}(0,\,\sigma_\text{step}^2)\;\hat{\mathbf{e}}$
      \Comment{Gaussian displacement in $x,y$}
    \State Wrap $\mathbf{r}_k$ with periodic boundaries (mod~$L$)
    \State $E_\text{new} \gets E_\text{mol}(\mathbf{r}_k, h)$
    \State $\Delta E \gets E_\text{new} - E_\text{old}$
    \If{$\Delta E \le 0$ \textbf{or}
        $\log\mathcal{U}(0,1) < -\Delta E$}
      \State Accept move
    \Else
      \State $\mathbf{r}_k \gets \mathbf{r}_\text{old}$
        \Comment{Reject}
    \EndIf
  \EndFor
  \Statex
  \Statex \Comment{\textbf{Phase 2: Membrane height updates}}
  \For{each grid cell $(i,j)$}
    \State $h_\text{old} \gets h_{ij}$
    \State $E_\text{mol,old} \gets$ sum of $E_\text{mol}$ for all
      molecules in cell $(i,j)$
    \State $h_{ij} \gets h_\text{old}
      + \mathcal{N}(0,\,\sigma_h^2)$;~~
      \textbf{if} $h_{ij}<0$ \textbf{then} $h_{ij}\gets -h_{ij}$
      \Comment{Reflect to keep $h\ge 0$}
    \State $\Delta E_\text{bend} \gets$
      local bending energy change (5-point stencil)
      \Comment{$O(1)$ via Eq.~\ref{eq:laplacian}}
    \State $E_\text{mol,new} \gets$ sum of $E_\text{mol}$ for all
      molecules in cell $(i,j)$ at new $h_{ij}$
    \State $\Delta E \gets \Delta E_\text{bend}
      + (E_\text{mol,new} - E_\text{mol,old})$
    \If{$\Delta E \le 0$ \textbf{or}
        $\log\mathcal{U}(0,1) < -\Delta E$}
      \State Accept move
    \Else
      \State $h_{ij} \gets h_\text{old}$ \Comment{Reject}
    \EndIf
  \EndFor
\EndFor
\Statex
\State \textbf{Output:} Depletion width
  $w = \text{median}(r^\text{CD45}) - \text{median}(r^\text{TCR})$
\end{algorithmic}
\end{algorithm}


\subsection{Step sizes}

Step sizes are derived from Brownian dynamics.  A stability-limited
time step $\Delta t$ is computed as
\begin{equation}
  \Delta t = \frac{\Delta x^2}{2\,D_h\,\kappa}\;\alpha_\text{safety}
\end{equation}
with safety factor $\alpha_\text{safety}=0.5$.  The molecular and
membrane height step sizes are then
\begin{equation}
  \sigma_\text{step} = \sqrt{2\,D_\text{mol}\,\Delta t},
  \qquad
  \sigma_h = \sqrt{2\,D_h\,\Delta t}
\end{equation}
where $D_\text{mol}=10^4$~nm$^2$/s and $D_h=5\times10^4$~nm$^2$/s
are the molecular and membrane height diffusion coefficients.
An alternative ``paper'' mode uses fixed $\Delta t = 0.01$~s and
$\sigma_h = 1$~nm.


\subsection{Efficient bending energy update}

Rather than recomputing the full bending energy (an $O(N_\text{grid}^2)$
operation), the change in bending energy when a single grid cell
$(i,j)$ is updated is computed in $O(1)$ time.  Only the five
Laplacian stencil cells affected by the change---$(i,j)$ itself and
its four periodic neighbours---need to be considered:

\begin{equation}
  \Delta E_\text{bend}
  = \frac{\kappa}{2}\,\Delta x^2
    \sum_{(a,b)\in\mathcal{N}(i,j)}
    \Bigl[\bigl(\nabla^2_d h_{ab}^\text{new}\bigr)^2
          - \bigl(\nabla^2_d h_{ab}^\text{old}\bigr)^2\Bigr]
\end{equation}

\noindent
where $\mathcal{N}(i,j) = \{(i,j),\,(i\pm1,j),\,(i,j\pm1)\}$ with
periodic wrapping.


\subsection{Checkerboard decomposition for Phase~2}

Phase~2 (membrane height updates) uses a red--black checkerboard
decomposition.  Grid cells are coloured so that each cell of a given
colour has only opposite-colour neighbours in the five-point stencil.
All same-colour cells can then be updated simultaneously without data
races, because their energy deltas depend only on the frozen
opposite-colour neighbours.

Each sweep processes both colours sequentially (red then black),
ensuring that every grid cell is updated exactly once.  On the CPU
this proceeds serially within each colour; on the GPU, all cells of a
given colour are dispatched in parallel.


\subsection{Initial conditions}

\begin{itemize}
  \item The membrane height field is initialised uniformly at
    $h = h_\text{init} = 70$~nm, with a central disc
    (radius $\sim N_\text{grid}/8$ cells) depressed to
    $h_0^\text{TCR}=13$~nm to seed a tight-contact zone.
  \item pMHC ligands are placed on the APC surface---either uniformly
    across the patch, or within a centred disc of radius $L/3$
    (``inner-circle'' mode, default).  They remain static throughout
    the simulation.
  \item TCR molecules are co-located with randomly chosen pMHC
    positions, with a small Gaussian jitter ($\sigma_\text{bind}$),
    so that each TCR starts near a cognate pMHC.  When no pMHC are
    present, TCRs are placed with a Gaussian bias towards the patch
    centre (spread $L/6$).
  \item CD45 molecules are placed uniformly at random across the
    patch.
\end{itemize}


\subsection{Output}

The primary observable is the \emph{depletion zone width}, defined as
the difference between the median radial distance from the patch
centre for CD45 molecules and that for TCR molecules:

\begin{equation}
  w_\text{depletion}
  = \max\!\Bigl(0,\;
    \widetilde{r}_\text{CD45} - \widetilde{r}_\text{TCR}\Bigr)
\end{equation}

\noindent
where $\widetilde{r}$ denotes the median over all molecules of that
species.  The use of the median rather than a percentile-based measure
provides robustness at small molecule counts.


%% ----------------------------------------------------------------
\section{Simulation parameters}
\label{sec:params}

Table~\ref{tab:params} lists all tunable parameters and their default
values from Supplementary Table~S1 of the reference.

\begin{table}[ht]
  \centering
  \caption{Default simulation parameters.}
  \label{tab:params}
  \begin{tabular}{@{}lll@{}}
    \toprule
    \textbf{Parameter} & \textbf{Symbol} & \textbf{Default value} \\
    \midrule
    Patch size          & $L$                   & 2000~nm ($2\;\mu$m) \\
    Grid resolution     & $N_\text{grid}$       & 64 \\
    Grid spacing        & $\Delta x$            & $L / N_\text{grid} \approx 31.25$~nm \\
    Number of TCR       & $N_\text{TCR}$        & 50 \\
    Number of CD45      & $N_\text{CD45}$       & 100 \\
    Number of pMHC      & $N_\text{pMHC}$       & 200 \\
    Association energy  & $U_\text{assoc}$      & $20\;k_BT$ \\
    Binding well width  & $\sigma_\text{bind}$  & 3~nm \\
    CD45 ectodomain     & $h_\text{CD45}$       & 50~nm \\
    Repulsive constant  & $k_\text{rep}$        & $1.0\;k_BT/\text{nm}^2$ \\
    Initial height      & $h_\text{init}$       & 70~nm \\
    TCR contact height  & $h_0^\text{TCR}$      & 13~nm \\
    Mol.\ diffusion     & $D_\text{mol}$        & $10^4$~nm$^2$/s \\
    Height diffusion    & $D_h$                 & $5\times10^4$~nm$^2$/s \\
    Safety factor       & $\alpha_\text{safety}$ & 0.5 \\
    Bending rigidity    & $\kappa$              & Swept (input parameter, $k_BT$) \\
    Simulation time     & $t$                   & Swept (input parameter) \\
    Random seed         & ---                   & 42 (per-point FNV-1a hash derived) \\
    \bottomrule
  \end{tabular}
\end{table}


%% ----------------------------------------------------------------
\section{Implementation}
\label{sec:implementation}

The model is implemented as a portable C99 simulation core with a
C++20 command-line interface and an optional GPU backend.  A single
binary (\texttt{ks\_gpu}) provides both CPU and GPU execution modes.
The build system uses CMake with platform detection: Apple Metal on
macOS, with a planned CUDA backend for Linux.

\subsection{Code architecture}

The source is organised into four layers (Table~\ref{tab:arch}).

\begin{table}[ht]
  \centering
  \caption{Source code organisation.}
  \label{tab:arch}
  \small
  \begin{tabular}{@{}llp{7cm}@{}}
    \toprule
    \textbf{Layer} & \textbf{Language} & \textbf{Files / purpose} \\
    \midrule
    Simulation core &
      C99 &
      \texttt{simulation.c}, \texttt{potentials.c}, \texttt{rng.c}
      --- MC sweep, molecular potentials, PCG64 RNG.
      Fully portable; no platform dependencies. \\[4pt]
    Shared physics &
      C99 / MSL &
      \texttt{ks\_physics.h} --- shared \texttt{static inline} float
      functions for TCR binding, CD45 repulsion, Laplacian, and
      bending-energy delta.  Preprocessor guards
      (\texttt{\_\_METAL\_VERSION\_\_}, \texttt{\_\_CUDA\_ARCH\_\_})
      select Metal address-space qualifiers or CUDA decorators. \\[4pt]
    GPU backend &
      Objective-C / MSL &
      \texttt{gpu\_engine.h} (backend-neutral C~API),
      \texttt{metal\_engine.m} (Metal implementation),
      \texttt{shaders.metal} (GPU compute kernels). \\[4pt]
    CLI &
      C++20 &
      \texttt{main.cpp} --- argument parsing, JSON I/O
      (nlohmann/json, header-only), directory management
      (\texttt{std::filesystem}), deterministic seed derivation
      (FNV-1a hash).  No platform-specific dependencies. \\
    \bottomrule
  \end{tabular}
\end{table}

The shared physics header (\texttt{ks\_physics.h}) is the single
source of truth for the four float-precision energy functions used by
both the CPU Phase~2 path (\texttt{simulation.c}) and the GPU
evaluate kernel (\texttt{shaders.metal}).  This eliminates the
previous code duplication and guarantees CPU--GPU equivalence by
construction.

\subsection{GPU backend interface}

The GPU backend is accessed through a four-function C~API defined in
\texttt{gpu\_engine.h}: \texttt{gpu\_engine\_create},
\texttt{gpu\_engine\_destroy}, \texttt{gpu\_engine\_h\_ptr}, and
\texttt{gpu\_engine\_grid\_update}.  The interface uses an opaque
\texttt{void*} handle, allowing different backends (Metal, CUDA) to be
linked at compile time without changing the simulation core.

\subsection{CPU mode}

In CPU mode, Phase~1 (molecular moves) and Phase~2 (grid height
updates) both execute sequentially.  Random numbers are generated with
the PCG64 algorithm, seeded deterministically from the user-provided
seed.  Phase~2 uses the checkerboard decomposition described above but
updates each colour's cells serially.

\subsection{GPU acceleration (Metal)}

On Apple Silicon, Phase~2 is offloaded to the GPU via Metal compute
shaders.  The pipeline consists of four kernel dispatches per colour:

\begin{enumerate}
  \item \textbf{Propose}: Each thread proposes a new height for one
    cell using a Philox counter-based RNG.
  \item \textbf{Snapshot}: Copies the full height field to a frozen
    buffer so that the evaluate kernel reads consistent neighbour
    values.
  \item \textbf{Evaluate}: Computes the energy delta (bending +
    molecular contributions) and the Metropolis acceptance criterion.
  \item \textbf{Apply}: Writes accepted proposals back to the
    authoritative height buffer.
\end{enumerate}

\noindent
Phase~1 remains on the CPU (few molecules, inherently sequential).
The height field resides in a Metal shared buffer
(\texttt{MTLResourceStorageModeShared}), allowing zero-copy access
from both CPU and GPU on unified memory.

A grid-substep parameter $K$ allows batching $K$ Phase~2 sweeps per
Phase~1 molecular move, reducing CPU--GPU synchronisation overhead.
All $K \times 8$ kernel dispatches are encoded into a single Metal
command buffer before committing.

\subsection{Performance}

Table~\ref{tab:perf} shows representative timings on Apple M3 for
500-step simulations.

\begin{table}[ht]
  \centering
  \caption{CPU vs GPU wall-clock time (500 steps, $\kappa=20$,
    $N_\text{TCR}=125$, $N_\text{CD45}=500$).}
  \label{tab:perf}
  \begin{tabular}{@{}rrrr@{}}
    \toprule
    $N_\text{grid}$ & CPU (s) & GPU (s) & Speedup \\
    \midrule
     64  & 0.14  & 0.06  & $2.3\times$ \\
    128  & 0.50  & 0.06  & $8.7\times$ \\
    256  & 1.78  & 0.07  & $25\times$ \\
    512  & 6.29  & 0.09  & $69\times$ \\
    \bottomrule
  \end{tabular}
\end{table}


%% ----------------------------------------------------------------
\bibliographystyle{plainnat}
\begin{thebibliography}{1}
\bibitem[Neve-Oz et~al.(2024)]{neveoz2024}
  Y.~Neve-Oz, E.~Sherman, and B.~Raveh.
  \newblock Bayesian metamodeling of early T-cell antigen receptor signaling
    accounts for its nanoscale activation patterns.
  \newblock \textit{Frontiers in Immunology}, 2024.
\end{thebibliography}

\end{document}
